\section{Step 7: software and first pulls}\label{sec:walk-software}
How the software is built and how to change it: Sect.~\ref{sec:software}.

\subsection{Download}
Both repositories into the \emph{same} folder (the server finds the UI next to it):
\begin{lstlisting}
mkdir IBMS_Offenburg && cd IBMS_Offenburg
git clone https://github.com/lore95/BumpemServer.git
git clone <BumpemUI repository URL>
\end{lstlisting}
\fillin{BumpemUI repository URL (not on GitHub yet)}
Large reference files are not in git (thesis, datasheets, photos, \file{legacy/}): copy them into
\file{BumpemServer/docs/thesis/}, \file{docs/reference/}, \file{docs/hardware-photos/}, \file{legacy/}.
\fillin{where they are stored}

\subsection{Install}
\begin{enumerate}
\item Install Miniconda (\url{https://docs.conda.io/en/latest/miniconda.html}). Linux only: install the Teensy
udev rule from pjrc.com.
\item In a terminal:
\begin{lstlisting}
conda create -n BumpemHome python=3.12
conda activate BumpemHome          # in every new terminal
cd IBMS_Offenburg/BumpemServer
pip install -e ".[server,dev]"
pip install platformio
pytest                             # 52 passed
\end{lstlisting}
\end{enumerate}

\subsection{Run the server locally, without hardware}
\begin{lstlisting}
bumpem serve --board sim
\end{lstlisting}
Open \url{http://127.0.0.1:8000/} in Chrome or Firefox. \emph{Connect}, \emph{Arm}, \emph{Perturb now}: the
live plot shows a pulse. \url{http://127.0.0.1:8000/docs} lists every API call. Stop the server with Ctrl+C.

\textbf{Check:} 52 tests pass and the simulated pulse appears in the plot.

\subsection{Flash the Teensy}
Ask the project lead first: flashing erases what is on the Teensy.
\begin{lstlisting}
bumpem ports                       # e.g. /dev/cu.usbmodem169222501 or COM7 = PORT
scripts/flash.sh                   # press the button on the Teensy if asked
bumpem info --port PORT            # proto=1, fw=1.0.0
bumpem monitor --port PORT         # about 100 lines per second, Ctrl+C to stop
\end{lstlisting}
If \cmd{info} says ``DAC init failed'': check DAC VCC, GND, SDA (pin 18), SCL (pin 19), replug USB. If
replies stop at random: try another USB port or cable before suspecting the software.

\subsection{Test pull per tower (testing mode)}
\begin{lstlisting}
bumpem serve --testing --board PORT --modules ABCD
\end{lstlisting}
Only list the towers that are plugged in (e.g.\ \cmd{--modules ACD}). Open \url{http://127.0.0.1:8000/}.
Order: USB, then 24\,V, then the tower's 48\,V. Nobody near the ropes. On the testing page: Pull 5\,N, then
10\,N, 400\,ms, one tower at a time.

\textbf{Check:} only that tower's rope moves; 10\,N pulls visibly harder than 5\,N.

\subsection{Closed loop (production mode)}
\begin{lstlisting}
bumpem serve --board PORT --modules ABCD
\end{lstlisting}
\begin{enumerate}
\item Connect. Ropes slack: \emph{Zero sensors}.
\item Settings: feedforward gain (\cmd{kff}) 2, baseline 2\,N, max force 40\,N. \emph{Save as default}
(the server writes them to the board at every connect).
\item Arm: each rope holds about 2\,N. \emph{Record}, perturbation 10\,N, stop recording: the CSV is in
\file{BumpemServer/Logs/}.
\item STOP (red button or Esc) and Release drop the rope tension: with a subject on the treadmill, stop the
treadmill first.
\end{enumerate}

\subsection{You have reached the current version when}
\begin{enumerate}
\item every tower passes the test pull;
\item at rest each force reads a few N and a hand pull changes only its own reading
(\file{scripts/force_check.py});
\item 6 / 8 / 12\,kg on each load cell read within about 2.5\,\% (Sect.~\ref{sec:calibration});
\item after a Teensy replug and Connect, the saved settings are back;
\item each tower holds 2\,N armed and a recorded 10\,N perturbation is in \file{Logs/};
\item the e-stop switches every ESCON off.
\end{enumerate}
State on 14 October 2026: item 1 passed on A, C and D; A and D held 2\,N. Items 2--4, 6 and recording on the
real board are still open; tower B waits for its new ESCON.
